\subsection{An actual rectangular-Kronecker quotient point with an exponential stabilizer}
\label{sec:tensor-chain-point}

\paragraph{P-F60 (the chain as a tensor, including its full product stabilizer).}
Let $n\ge4$, $V=\mathbb C^n$, $W=V^{\otimes3}$,
$G=SL_n^3$, and $M=2^{n-1}$.  Set
\[
 r_0=0,\qquad r_i=\frac{1-r_{i-1}}2\quad(1\le i<n),
 \qquad r_n=1-r_{n-1},\qquad c_i=\sqrt{r_i}>0,
                                                        \tag{P-F60a}
\]
and define the symmetric tensor
\[
 T_n=\sum_{i=1}^{n-1}c_i
   (e_i\otimes e_i\otimes e_{i+1}
    +e_i\otimes e_{i+1}\otimes e_i
    +e_{i+1}\otimes e_i\otimes e_i)
             +c_n e_n^{\otimes3}.                      \tag{P-F60b}
\]
Its $G$-orbit is closed and nonzero.  Put
$w_i=(-2)^{i-1}$ for $1\le i<n$, $w_n=0$, and
$D_\zeta=\operatorname{diag}(\zeta^{w_1},\ldots,\zeta^{w_n})$.
Its full stabilizer in $GL_n^3$ is
\[
 \operatorname{Stab}_{GL_n^3}(T_n)
   =\{(aD_\zeta,bD_\zeta,cD_\zeta):
          a,b,c\in\mathbb C^*,\ abc=1,\ \zeta^M=1\}.
                                                        \tag{P-F60c}
\]
In particular this formula includes all independent scalar
factors and all disconnected components; it is not just
a calculation inside a diagonal subgroup.  For
$\chi(A,B,C)=\det A\det B\det C$,
\[
 \chi\bigl(\operatorname{Stab}_{GL_n^3}(T_n)\bigr)
                         =\boldsymbol\mu_M,
 \qquad
 \operatorname{Stab}_{G}(T_n)
   =K_n:=\{(aI,bI,cI):a^n=b^n=c^n=1,\ abc=1\}.
                                                        \tag{P-F60d}
\]
The group $K_n$ is precisely the kernel of $G$ on $W$.

Let
\[
 R^{\mathrm{ten}}_n=\mathbb C[W]^G
       =\bigoplus_{s\ge0}(R^{\mathrm{ten}}_n)_s,
 \qquad
 (R^{\mathrm{ten}}_n)_s=\mathbb C[W]^G_{\text{ordinary degree }ns}.
                                                        \tag{P-F60e}
\]
All nonzero ordinary degrees are divisible by $n$, and
\[
 \dim(R^{\mathrm{ten}}_n)_s=k((s^n),(s^n),(s^n)).        \tag{P-F60f}
\]
If $X_n=\operatorname{Spec}R^{\mathrm{ten}}_n$ and
$p_n$ is the image of $T_n$, then $p_n\ne0$, $X_n$ is
smooth at $p_n$, and its grading stabilizer is exactly
$\boldsymbol\mu_M$.  The fixed scheme
$X_n^{\boldsymbol\mu_M}$ is smooth of dimension one at
$p_n$.  The closure $Y_n$ of its nonzero grading orbit
is consequently an irreducible curve component of this
fixed scheme.  Its regraded leading Hilbert coefficient
is $e_{Y_n}=1$.

These are assertions about an actual Kronecker invariant
ring and an actual component of its fixed scheme.
The multiplication-compatible canonical-ring identification
below also proves that every homogeneous system of parameters
has a degree divisible by $M$.  In particular the canonical
Kronecker modular index is at least $2^{n-1}$.
They do not yet assert a pole of the global Hilbert
series: other components have not been classified, and
global cancellation has not been excluded.

\paragraph{Balance and closedness.}
The recurrence has the explicit solution
$r_i=(1-(-1/2)^i)/3$ for $i<n$.
Hence $1/4\le r_i\le1/2$ and $r_n\ge1/2$.
No two distinct ordered triples in the support of $T_n$
agree in two positions.  Indeed, a pair of equal indices
$(i,i)$ extends only to the remaining index $i+1$
when $i<n$, or to $n$ when $i=n$; a pair of distinct
indices occurs only when it is $(i,i+1)$ in either order,
and then its remaining index is $i$.  Thus all three
unnormalized reduced density matrices of $T_n$ are
diagonal.  Their diagonal entries are, respectively,
\[
 2r_1=1,\qquad r_{i-1}+2r_i=1\ (2\le i<n),
                     \qquad r_{n-1}+r_n=1.
                                                        \tag{P-F60g}
\]
They are therefore all $I_n$, and $\|T_n\|^2=n$.
For any traceless Hermitian matrices $H_1,H_2,H_3$,
\[
 \left\langle
 (H_1\otimes I\otimes I+I\otimes H_2\otimes I
                         +I\otimes I\otimes H_3)T_n,T_n
 \right\rangle
                  =\operatorname{tr}H_1+\operatorname{tr}H_2
                                      +\operatorname{tr}H_3=0.
                                                        \tag{P-F60h}
\]
This is exactly vanishing of the $SU_n^3$ moment map,
with the product tensor Hermitian norm.  The Kempf--Ness
criterion implies that the complex $SL_n^3$ orbit is
closed.  A primary proof with this general representation
scope is Losev, \emph{The Kempf--Ness theorem and invariant
theory}, Theorem~4, printed p.~1 and its proof on p.~2,
\url{https://arxiv.org/pdf/math/0605756}.
We use the theorem on the explicitly computed moment map;
torus semistability alone would not prove closedness here.

\paragraph{The centroid is scalar.}
For a symmetric tensor $T$ define its centroid by
\[
 C(T)=\{L\in\operatorname{End}(V):
          (L\otimes I\otimes I)T=(I\otimes L\otimes I)T\}.
\]
The matrix slices of $T_n$ in the third tensor position are
\[
 \begin{split}
 S_1&=c_1(E_{12}+E_{21}),\\
 S_k&=c_{k-1}E_{k-1,k-1}
                    +c_k(E_{k,k+1}+E_{k+1,k})\quad(2\le k<n),\\
 S_n&=c_{n-1}E_{n-1,n-1}+c_nE_{nn}.
 \end{split}                                         \tag{P-F60i}
\]
The centroid condition is that $LS_k=S_kL^t$, or
equivalently that each $LS_k$ is symmetric.
Each $S_k$ is invertible on its coordinate support,
which is $\{1,2\}$, $\{k-1,k,k+1\}$, or
$\{n-1,n\}$, respectively.  Since $LS_k$ has zero
columns outside this support, symmetry implies zero
rows there too.  Invertibility on the support then
implies that every column of $L$ whose index belongs
to the support has no entries outside the support.
Intersect these restrictions for all slices containing
that column.  For columns $2,\ldots,n-1$ the intersection
is the singleton consisting of the column index.
For column~1 it is $\{1,2\}$, and for column~$n$ it
is $\{n-1,n\}$.  This assertion also holds for $n=4$,
where the two interior slice supports overlap in two
coordinates but their intersections with the end
supports still leave singletons.

Thus the only possible off-diagonal entries are
$L_{21}$ and $L_{n-1,n}$.  Symmetry of $LS_2$ in
positions $(2,1)$ and $(1,2)$ gives
$c_1L_{21}=c_2L_{13}=0$.  Symmetry of $LS_n$ in
positions $(n-1,n)$ and $(n,n-1)$ gives
$c_nL_{n-1,n}=c_{n-1}L_{n,n-1}=0$.
All $c_i$ are nonzero, so $L$ is diagonal.
Finally, the $(k,k+1)$ and $(k+1,k)$ entries of
$LS_k$, for $1\le k<n$, force $L_{kk}=L_{k+1,k+1}$.
We have proved
\[
                         C(T_n)=\mathbb C I.            \tag{P-F60j}
\]

\paragraph{From the centroid to the full product stabilizer.}
Suppose $(A,B,C)T_n=T_n$.  Symmetry of $T_n$ implies
$(B,A,C)T_n=T_n$ as well.  Apply $(A^{-1},B^{-1},C^{-1})$
to obtain
$(L,L^{-1},I)T_n=T_n$, where $L=A^{-1}B$.
Multiplication in the second factor by $L$ gives
$(L,I,I)T_n=(I,L,I)T_n$, so (P-F60j) implies
$B=bA$ for some $b\ne0$.  Swapping the second and third
positions in the same way implies $C=cA$.
It follows that $A^{\otimes3}T_n=(bc)^{-1}T_n$.
Thus the projective class of $A^t$ is an automorphism,
in the coordinate-substitution convention, of the symmetric cubic
\[
                   3\sum_{i<n}c_i x_i^2x_{i+1}+c_nx_n^3.
                                                        \tag{P-F60k}
\]
A diagonal change of variables takes this cubic to the
chain $F_n$ of P-F24: solve first for the last variable
using the nonzero coefficient $c_n$, and then solve
successively backwards for the square coefficient of
each earlier variable.  All coefficients are nonzero,
so this gives an invertible diagonal change.  The full
projective automorphism group of (P-F60k) is therefore
the same diagonal cyclic group found in P-F24, since
diagonal conjugation leaves that group unchanged.
Its representatives are $D_\zeta$, $\zeta^M=1$;
these are diagonal, so transposition makes no change.
Consequently $A=aD_\zeta$, $B=b'D_\zeta$,
$C=c'D_\zeta$ for independent nonzero scalars.

For every $i<n-1$, $2w_i+w_{i+1}=0$; also
$2w_{n-1}=\pm M$ and $w_n=0$.  Hence
$D_\zeta^{\otimes3}T_n=T_n$.  The condition on the
three scalar factors is exactly $ab'c'=1$.
This proves both inclusions in (P-F60c), with no
connectedness assumption on the original stabilizer.
If
\[
                S=\sum_{i=1}^{n}w_i,
           \qquad 3S=1-(-2)^{n-1}\equiv1\pmod M,
\]
then $\chi(aD_\zeta,bD_\zeta,cD_\zeta)=\zeta^{3S}=\zeta$.
This proves the first assertion in (P-F60d).
An element in $SL_n^3$ has $\chi=1$, so necessarily
$\zeta=1$ and the three scalars have $n$th power one.
This proves the second assertion.  Conversely these
scalar triples act trivially on every tensor, and any
triple acting trivially on $V^{\otimes3}$ has scalar
factors with product one: fix two nonzero pure factors
and vary the third to see that each factor preserves
every line, hence is scalar.  Thus $K_n$ is the full
ineffective kernel.

The character $\chi$, the product action, and its
stabilizer period agree exactly with B\"urgisser--Ikenmeyer,
\emph{Fundamental invariants of orbit closures},
equations~(4.1)--(4.2) and Definition~4.1, printed
pp.~18--19, \url{https://arxiv.org/pdf/1511.02927v2}.
In particular, the exponential order here is a product
determinant character, not the order of an arbitrarily
chosen matrix lift.

\paragraph{The quotient, its grading, and the fixed tangent space.}
The scalar center of any one $SL_n$ factor forces every
nonzero homogeneous invariant degree to be divisible
by $n$.  In degree $ns$, the triple Schur--Weyl
decomposition of $\operatorname{Sym}^{ns}(V^{\otimes3})$
has an $SL_n^3$-invariant summand precisely when its
three Schur partitions are all $(s^n)$.  Each corresponding
determinant module is one-dimensional.  Dualizing to
coordinate functions preserves the dimension of invariants.
This proves (P-F60e)--(P-F60f).  The same identification
is stated in B\"urgisser--Ikenmeyer, equation~(5.2),
printed p.~23.

Because the $G$-orbit of $T_n$ is nonzero and closed,
it is distinct from the closed orbit at the origin;
uniqueness of the closed orbit in a reductive quotient
fiber gives $p_n\ne0$.  For $A=(A_1,A_2,A_3)\in GL_n^3$
and $F\in(R^{\mathrm{ten}}_n)_s$, write
$A_i=\alpha_i B_i$ with $B_i\in SL_n$.
The point-evaluation convention gives
\[
                   F(AT)=\chi(A)^s F(T).                \tag{P-F60l}
\]
It is independent of the choices of $n$th roots
$\alpha_i$, since $\deg F=ns$.  Thus the action induced
on $X_n$ is the grading action with parameter $\chi(A)$.
Formula (P-F60d) shows that $\boldsymbol\mu_M$ fixes
$p_n$.  Conversely, if $u\in\mathbb C^*$ fixes $p_n$,
choose $A\in GL_n^3$ with $\chi(A)=u$.
The tensor $AT_n$ has the same quotient image as $T_n$.
Both $G$-orbits are closed, since $GL_n^3$ normalizes $G$.
They therefore coincide.  After multiplying $A$ by an
element of $G$, it belongs to the stabilizer of $T_n$,
and (P-F60d) gives $u\in\boldsymbol\mu_M$.  This proves
the exact grading stabilizer.

Pass to the effective reductive group $G/K_n$.
Its stabilizer at $T_n$ is trivial.  Luna's slice theorem
therefore identifies an \'etale neighborhood of $p_n$
with a smooth normal slice, and gives the tangent
exact sequence
\[
 0\longrightarrow\mathfrak{sl}_n^3 T_n
     \longrightarrow W\longrightarrow T_{p_n}X_n
       \longrightarrow0.                              \tag{P-F60m}
\]
The orbit map on the left is injective because the
stabilizer has zero Lie algebra.  In particular,
$\dim X_n=n^3-3n^2+3$ at $p_n$.
The precise slice locator is Luna, \emph{Slices \'etales},
Section~III.1, the theorem on p.~97 and Remark~$1^\circ$
on pp.~97--98; the unique-closed-orbit assertion used
above is Section~I.1, p.~84.
The primary text is
\url{https://www.numdam.org/item/10.24033/msmf.110.pdf}.

Take a primitive $M$th root $\zeta$.
The triple $(D_\zeta,D_\zeta,D_\zeta)$ fixes $T_n$,
normalizes $G$, and induces the grading generator
on $X_n$ by (P-F60l).  Its eigenvalue-one tensors in
$W$ are exactly the $3(n-1)$ ordered chain terms and
$e_n^{\otimes3}$.  Indeed the nonzero weights $w_i$
have distinct $2$-adic valuations.  In a vanishing
sum of three weights, the least valuation cannot occur
once; it must occur twice, and then the remaining
weight is its negative double.  This gives the chain
terms, including the last nonzero weight followed by
zero.  Three equal nonzero weights cannot sum to zero,
and three zero weights give the last term only.
All eigenvalues of $D_\zeta$ on $V$ are distinct, so
the fixed subspace of each adjoint $\mathfrak{sl}_n$
is the traceless diagonal subspace of dimension $n-1$.
Taking fixed vectors in (P-F60m) is exact in
characteristic zero, and yields
\[
           \dim(T_{p_n}X_n)^{\boldsymbol\mu_M}
                        =(3n-2)-3(n-1)=1.              \tag{P-F60n}
\]
The fixed scheme of a finite-order automorphism of a
smooth complex variety is smooth at a fixed point:
average local parameters into eigenspaces, and its
completed fixed ideal is generated by the parameters
with nontrivial eigenvalues.  Thus the fixed scheme
has dimension one at $p_n$.  The grading orbit has
dimension one, since its stabilizer is finite; it is
contained in this fixed scheme, and its irreducible
closure is consequently a curve component $Y_n$.
Every point of the nonzero orbit remains a smooth
point of $X_n$ by grading translation.

\paragraph{Leading density and the remaining pole hypothesis.}
Restriction to the grading orbit maps a homogeneous
invariant $F_s$ to $F_s(p_n)u^s$ in
$\mathbb C[u,u^{-1}]$.  Its image is the semigroup
algebra $\mathbb C[u^s:s\in S_n]$, where
$S_n=\{s\ge0:\exists F_s,\ F_s(p_n)\ne0\}$.
A finite set of homogeneous algebra generators for
$R^{\mathrm{ten}}_n$ shows that $S_n$ is generated by
the degrees of those generators nonzero at $p_n$:
every nonzero evaluated homogeneous polynomial has
at least one nonzero evaluated monomial, and products
of nonzero evaluations stay nonzero.  The stabilizer
of this grading orbit is $\boldsymbol\mu_{\gcd(S_n)}$.
The exact stabilizer already proved implies
$\gcd(S_n)=M$.

The finitely generated semigroup $S_n/M$ of
nonnegative integers has gcd one and is cofinite.
To see this directly, take a positive element $h$.
Its generator residues generate the full finite cyclic
group modulo $h$, and their additive monoid is also
that group.  Choose one semigroup representative in
each residue class; all larger integers in that class
are obtained by adding $h$.  Therefore the Hilbert
function of $Y_n$ is eventually one at multiples of
$M$ and zero elsewhere.  In particular
\[
              \lim_{t\to1}(1-t^M)H_{Y_n}(t)=1,
                           \qquad e_{Y_n}=1.            \tag{P-F60o}
\]
We do not assert that $Y_n$ is normal at the origin
or that its semigroup has no initial holes.

At a primitive $M$th root the nonzero conormal
determinant of this curve component can be written
explicitly, with the point-action convention of
(P-F60l), as
\[
 C^{\mathrm{loc}}_n(\zeta)=
 \frac{\displaystyle\prod_{i\ne j}
            (1-\zeta^{w_i-w_j})^3}
      {\displaystyle\prod_{i,j,k:\ w_i+w_j+w_k\not\equiv0\ (M)}
            (1-\zeta^{w_i+w_j+w_k})}.
                                                        \tag{P-F60p}
\]
Here the ring automorphism is the pullback $F\mapsto
F\circ(D_\zeta,D_\zeta,D_\zeta)$, so on degree $s$
it has eigenvalue $\zeta^s$ by (P-F60l), exactly as
in P-F23.  On the local cotangent space this pullback
is the transpose of the point differential and has
the same eigenvalues, not their inverses.
Remove the fixed tangent line from (P-F60m) and use
multiplicativity of this moving determinant to obtain
(P-F60p).
Every displayed factor is nonzero.  If $Y_n$ were
the unique positive-dimensional component of
$X_n^{\boldsymbol\mu_M}$, P-F23 and (P-F60o) would
give an actual simple Hilbert pole with normalized
coefficient $C^{\mathrm{loc}}_n(\zeta)$.
That additional global hypothesis is refuted by P-F64 below,
which constructs a second positive-dimensional fixed component.
In particular (P-F60p) is a local conormal factor,
not an asserted residue of the full Hilbert series;
the numerical positive-form clause of ordinary Kronecker
PH3 and its strict positivity/saturation indices are not
concluded from this local calculation.

\paragraph{The actual canonical Kronecker algebra and its parameter degrees.}
The stabilizer already proved has a consequence that does
not require a global pole.  Every homogeneous system of
parameters of $R^{\mathrm{ten}}_n$ has at least one degree
divisible by $M$.  To apply this to the source's canonical
modular index, we first identify the graded algebra,
including its multiplication.

Put $E=V^*$ and $U=E\otimes E$.  Thus the coordinate ring
of our original tensor space is
$\mathbb C[W]=\operatorname{Sym}(U\otimes E)$.
The contragredient images of the original three $SL(V)$
factors are the full three groups $SL(E)$; no equivariant
identification of $V$ with $V^*$ is being asserted.
Cauchy's decomposition gives
\[
 \operatorname{Sym}^{d}(U\otimes E)
     =\bigoplus_{\alpha\vdash d}
           \mathbb S_\alpha U\otimes\mathbb S_\alpha E.
\]
Taking invariants under the third $SL(E)$ retains precisely
$d=ns$ and $\alpha=(s^n)$.  Choose a nonzero volume vector
$\omega\in\det(E)$ and use its powers to trivialize the
one-dimensional factors.  We obtain
\[
 \operatorname{Sym}(U\otimes E)^{SL(E)_{\rm third}}
       \simeq \bigoplus_{s\ge0}\mathbb S_{(s^n)}U,
 \qquad \text{ordinary degree }ns\text{ on the left has degree }s.
                                                        \tag{P-F60q}
\]
This is a graded algebra isomorphism with the Cartan
multiplication on the right, not just a character identity.
Indeed, multiplication from degrees $s,t$ to $s+t$ is
$GL(U)$-equivariant.  The Cartan summand of highest weight
$((s+t)^n)$ occurs once in
$\mathbb S_{(s^n)}U\otimes\mathbb S_{(t^n)}U$.
The product of two nonzero highest vectors is nonzero in
the polynomial domain and is again a highest vector.
Choose a highest vector $z$ in degree one and normalize
the highest vector in degree $s$ as $z^s$; these products
are nonzero and provide compatible normalizations for
every $s,t$.  Consequently each multiplication is exactly
the normalized Cartan projection.  Removing the determinant
factors by $\omega^s\omega^t=\omega^{s+t}$ respects the
same multiplication.

Let $H=GL(E)\times GL(E)$ act on $U$ and
$G'=GL(U)=GL_{n^2}$.  The source Kronecker input is
\[
       \lambda=(1^n),\qquad
       \pi=((1^n),(1^n)),\qquad
       \rho:H\longrightarrow G',\quad(A,B)\longmapsto A\otimes B.
                                                        \tag{P-F60r}
\]
Every vector of $\mathbb S_{(s^n)}U$ has scalar weight
$ns$ under each of the two $GL(E)$ factors.
An $SL(E)\times SL(E)$-invariant vector therefore
transforms under $H$ by $\det^s\boxtimes\det^s$.
This follows either from the one-dimensional characters
of $GL_n/SL_n$, or by writing each invertible matrix as
a scalar times an $SL_n$ matrix; the scalar ambiguity
is removed by the divisibility $n\mid ns$.
Taking the two remaining invariants in (P-F60q) therefore
identifies our ring, with its multiplication, as
\[
 R^{\mathrm{ten}}_n\simeq
 \bigoplus_{s\ge0}
 \left(\mathbb S_{(s^n)}U\otimes
       \bigl((\det E)^{-s}\boxtimes(\det E)^{-s}\bigr)\right)^H
       =T(m_\lambda^\pi).
                                                        \tag{P-F60s}
\]
Here determinant-line bases are multiplied by their powers,
so the displayed identification also respects the Cartan
product on the target-character factors.
This is exactly the canonical algebra in GCT VI,
\S4.1.2, equation~(4.5), printed p.~69; the Kronecker
embedding is specified in Problem~1.1.1.  Its canonical
modular index is defined from a lexicographically minimal
h.s.o.p. of this $T$ in \S4.1.1, printed p.~68:
\url{https://arxiv.org/pdf/0704.0229}.

Let $h_1,\ldots,h_d$ be any homogeneous system of
parameters of this positively graded algebra.  The quotient
$R^{\mathrm{ten}}_n/(h_1,\ldots,h_d)$ is finite-dimensional.
Every positive-degree element of that quotient is nilpotent,
so the common zero set of the parameters is just the
vertex of $X_n$.  Since $p_n\ne0$, some $h_i(p_n)$ is
nonzero.  For every $\zeta\in\boldsymbol\mu_M$, the exact
grading stabilizer of $p_n$ gives
\[
 h_i(p_n)=h_i(\zeta\cdot p_n)
                 =\zeta^{\deg h_i}h_i(p_n).
\]
Thus $M\mid\deg h_i$.  In particular this holds for the
source's lexicographically minimal parameter sequence,
and hence
\[
 \chi\bigl(k((1^n),(1^n),(1^n))\bigr)\ge2^{n-1}.
                                                        \tag{P-F60t}
\]
The notation in (P-F60t) denotes the canonical modular
index attached to that input and its stretching algebra;
it does not assert that the initial coefficient is nonzero.

The exact unified statement is GCT VI Hypothesis~3.3.4,
printed p.~51.  Its second, more-specific clause requires
this canonical index to be polynomial in $\operatorname{rank}(x)$.
The definition in \S3.4, printed p.~52, sums the ranks
of $H,G'$ and the two weight heights.  For (P-F60r) it is
\[
             \operatorname{rank}(x)=2n+n^2+n+2n=n^2+5n.
                                                        \tag{P-F60u}
\]
All parts of the three rectangular input partitions are
one, so their total partition bit length is also linear
in $n$.  The lower bound (P-F60t) consequently refutes
that canonical clause on an actual triple-rectangular
Kronecker family.  Even the larger introductory parameter
$\dim G'+O(n)=n^4+O(n)$ is polynomial in $n$.
This argument does not refute the first clause, which
allows an arbitrary positive rational form with small
denominator weights.  The period-collapse examples
P-F21--P-F22 show why a canonical parameter lower bound
alone cannot establish that numerical conclusion.
The global fixed-component and noncancellation problem
above, and any strict Kronecker index lower bound,
remain open here.

\paragraph{The canonical obstruction also occurs at positive initial inputs.}
For each $n\ge4$ there exists an odd positive integer
$t=t(n)$ such that
\[
        k((t^n),(t^n),(t^n))>0.
                                                        \tag{P-F60v}
\]
The required known result is B\"urgisser--Ikenmeyer,
\emph{Fundamental invariants of orbit closures},
Theorem~4.2, printed p.~19, Proposition~4.10, p.~22,
and Theorem~5.3 together with the discussion following
equation~(5.1), p.~23:
\url{https://arxiv.org/pdf/1511.02927}.
They prove that generic cubic-format tensors have
stabilizer period one and are polystable, and that the
ordinary invariant degree monoid generates $n\mathbb Z$.
After our regrading this says that
$\{s:(R^{\mathrm{ten}}_n)_s\ne0\}$ has gcd one.
If every positive member were even, that gcd would be
at least two.  This proves (P-F60v).

Fix such an odd $t$ and form the Veronese algebra
\[
       R^{[t]}_n=\bigoplus_{s\ge0}(R^{\mathrm{ten}}_n)_{ts},
       \qquad (R^{\mathrm{ten}}_n)_{ts}
                    \text{ has new degree }s.
                                                        \tag{P-F60w}
\]
The multiplication proof in (P-F60q)--(P-F60s) identifies
this with the actual canonical Kronecker algebra of
the input $((t^n),(t^n),(t^n))$.
Let $p_n^{[t]}$ be the image of $p_n$ under the finite
Veronese quotient map.  It is nonzero: a positive-degree
homogeneous $h$ nonzero at $p_n$ has $h^t\in R^{[t]}_n$
and $h^t(p_n)\ne0$.
If an element of new degree $s$ is nonzero at this point,
its original degree $ts$ must be divisible by $M$.
Since $t$ is odd and $M$ is a power of two, $M\mid s$.
In fact the grading stabilizer is still exactly
$\boldsymbol\mu_M$: (P-F60o)'s semigroup argument shows
that all sufficiently large original multiples of $M$
occur nonvanishingly at $p_n$, and hence all sufficiently
large new multiples of $M$ do so at $p_n^{[t]}$.
The same parameter-ideal argument now proves
\[
      \chi\bigl(k((t^n),(t^n),(t^n))\bigr)\ge2^{n-1}.
                                                        \tag{P-F60x}
\]
Moreover a nonzero invariant in original degree $t$
has nonzero powers, so
\[
 k(((ts)^n),((ts)^n),((ts)^n))>0\qquad(s\ge1).
\]
Thus the canonical obstruction persists even when every
realized positive stretch is nonzero.  All partition
heights and $\operatorname{rank}(x)=n^2+5n$ are unchanged.
No bound on the size of this chosen odd $t$ has been
proved here, so this refinement makes no claim about
an input-bit-length lower bound or an effective choice
of $t$.  It still gives no numerical positive-form or
strict-index conclusion.
